%------------------------------------------------------------------------------%
\begin{question}
  Calcule
  \(
    \int \frac{3x + 5}{(x+1)(x+2)} \, dx
  \)
\end{question}

%------------------------------------------------------------------------------%
\begin{resolution}{Solução}
  O grau do numerador
  \[
    P(x) = 3x + 5
  \]
  \pause
  é menor do que o do denominador
  \[
    Q(x) = (x+1)(x+2)
  \]
  \pause
  O denominador já está fatorado

  \pause
  Por frações parciais temos
  \[
    \pause
    f(x)
    \;=\; \frac{P(x)}{Q(x)}
    \pause
    \;=\; \frac{3x + 5}{(x+1)(x+2)}
    \pause
    \;=\; \frac{A}{x+1} + \frac{B}{x+2}
  \]
\end{resolution}

%------------------------------------------------------------------------------%
\begin{resolution}{Solução}
  \begin{align*}
    \onslide<+->{\frac{3x + 5}{(x+1)(x+2)} & =  \frac{A}{x+1} + \frac{B}{x+2}      \\[3mm]}
    \onslide<+->{3x + 5 & = \left(\frac{A}{x+1} + \frac{B}{x+2}\right)(x+1)(x+2)   \\[1mm]}
    \onslide<+->{& = A(x+2) + B(x+1)                                               \\[1mm]}
    \onslide<+->{& = Ax + 2A + Bx + B                                              \\[1mm]}
    \onslide<+->{& = (A+B)x + (2A + B)}
  \end{align*}
\end{resolution}

%------------------------------------------------------------------------------%
\begin{resolution}{Solução}
\begin{columns}
\column{0.5\linewidth}
  \onslide<+->{Para que os polinômios sejam iguais
  \[
    3x + 5 = (A+B)x + (2A + B)
  \]}
  \onslide<+->{os coeficientes precisam ser iguais}
  \[
    \onslide<+->{
    \systeme[sort={A,B}]{
      A + B = 3,
      2A + B = 5
    }}
  \]
\column{0.5\linewidth}
  \[
    \onslide<+->{B = 3 - A}
  \]
  \vspace{-1.5\baselineskip}
  \begin{align*}
    \onslide<+->{2A + B     & = 5     \\}
    \onslide<+->{2A + 3 - A & = 5     \\}
    \onslide<+->{A          & = 5 - 3}
    \onslide<+->{             = 2}
  \end{align*}
  \vspace{-1.5\baselineskip}
  \begin{align*}
    \onslide<+->{B & = 3 - A \\}
    \onslide<+->{  & = 3 - 2   }
    \onslide<+->{    = 1}
  \end{align*}
\end{columns}
\end{resolution}

%------------------------------------------------------------------------------%
\begin{resolution}{Solução}
\begin{align*}
  \onslide<+->{f(x)}
  \onslide<+->{& = \frac{3x + 5}{(x+1)(x+2)}           \\[2mm]}
  \onslide<+->{& = \frac{A}{x+1} + \frac{B}{x+2}       \\[2mm]}
  \onslide<+->{& = \frac{2}{x+1} + \frac{1}{x+2}}
\end{align*}
\end{resolution}

%------------------------------------------------------------------------------%
\begin{resolution}{Solução}
  \begin{align*}
    \onslide<+->{F(x)}
    \onslide<+->{& = \int \frac{3x + 5}{(x+1)(x+2)} dx              \\[2mm]}
    \onslide<+->{& = \int \frac{2}{x+1} + \frac{1}{x+2} dx          \\[2mm]}
    \onslide<+->{& = 2\int \frac{1}{x+1} dx + \int \frac{1}{x+2} dx \\[2mm]}
    \onslide<+->{& = 2\ln\abs{x+1} + \ln\abs{x+2} + c               \\[2mm]}
    \onslide<+->{& = \color{gray} \ln\left((x+1)^2\,\abs{x+2}\right) + c}
  \end{align*}
\end{resolution}

%------------------------------------------------------------------------------%
